\documentclass[11pt,letterpaper]{article}

\usepackage[utf8]{inputenc}
\usepackage[T1]{fontenc}
% Removed mathptmx to use LaTeX standard Computer Modern font like Jiaqi Chen
\usepackage[left=1in,right=1in,top=1in,bottom=1in]{geometry}
\usepackage{enumitem}
\usepackage{titlesec}
\usepackage{hyperref}
\usepackage{xcolor}

% Spacing and formatting
% Removed setstretch to use default compact line spacing
\usepackage[none]{hyphenat}
\sloppy

% Color setup for links
\hypersetup{
    colorlinks=true,
    linkcolor=black,
    filecolor=black,      
    urlcolor=blue,
    citecolor=black,
}

% Section formatting to match Jiaqi Chen (Small Caps)
\titleformat{\section}{\large\scshape}{}{0em}{}[\titlerule]
\titlespacing*{\section}{0pt}{2.5ex}{0.5ex}

% Custom commands for entries
\newcommand{\cventry}[4]{
  \noindent\textbf{#1} \hfill #2 \\
  \noindent\textit{#3} \hfill \textit{#4}
}

\newcommand{\cvproject}[3]{
  \noindent\textbf{#1} $|$ \textit{#2} \hfill #3
}

\newcommand{\cvpub}[1]{
  \noindent #1
}

% Remove page numbers
\pagestyle{empty}

\begin{document}

% HEADER
\begin{center}
    {\LARGE \textbf{THANH TRINH-XUAN}} \\[1ex]
    Viet Nam $|$ +84 867 663 236 $|$ \href{mailto:txthanh1005@gmail.com}{txthanh1005@gmail.com} $|$ \href{https://github.com/txthanh1005-maker/CV_design/tree/main/Research_Portfolio}{Thanh Trinh-Xuan Github}
\end{center}

\vspace{1ex}

% EDUCATION
\section*{Education}

\cventry{Hanoi University of Science and Technology (HUST)}{Hanoi, Vietnam}{Bachelor of Engineering in Electrical Engineering}{Aug. 2026 (Expected)}
\begin{itemize}[leftmargin=0.25in, itemsep=0pt, topsep=0pt]
    \item \textbf{GPA:} 3.2/4.0
    \item \textbf{Research interests:} mathematical programming and distributed algorithms in power systems, peer-to-peer energy markets in multi-microgrids, optimization and operation control in resilience-oriented interconnected smart grids
    \item \textbf{Affiliation:} Member of 100 Renewable Energy Laboratory.
    \item \textbf{Bachelor's Thesis:} Resilience-Oriented Peer-to-Peer Energy Trading in Multi-Microgrids via Network-Constrained Distributed Predictive Control
    \item \textbf{Relevant Coursework:} Power System Analysis, Power System Automation, Power System Protection and Control, Digital System Design, Programming Techniques.
\end{itemize}

\vspace{1ex}

\cventry{Lam Son High School for the Gifted Students}{Thanh Hoa, Vietnam}{High School Diploma, Mathematics Specialization}{Jul. 2022}

\vspace{1ex}

% RESEARCH AND WORK EXPERIENCE
\section*{Research Experience}
\noindent {Undergraduate Researcher} \hfill \textit{Sep. 2024 -- Present}
% \cvproject{Undergraduate Researcher}{Sep. 2024 -- Present}

\noindent\textit{Power Grid and Renewable Energy Laboratory, Hanoi, Vietnam}\\
\noindent\textit{Supervisor: Assoc. Prof. Nguyen Duc Tuyen}

\vspace{1ex}
\noindent \parbox[t]{0.82\linewidth}{\textbf{Bachelor's Thesis}}%
\hfill\parbox[t]{0.16\linewidth}{\raggedleft\textit{2025 -- 2026}}
\begin{itemize}[leftmargin=0.25in, itemsep=0pt, topsep=0pt]
    \item Modeled P-Q-V dynamics using \textbf{Second-Order Cone Programming (SOCP)} and integrated load shedding penalty functions to mitigate relaxation inexactness (fictitious currents), ensuring exact grid constraint enforcement.
    \item Implemented \textbf{Analytical Target Cascading (ATC)} to maintain decentralized coordination when extreme faults cause network fragmentation and render centralized control unviable.
    \item Designed \textbf{Model Predictive Control (MPC)} frameworks for dynamic rescheduling to mitigate immediate power imbalances during blackouts that invalidate static day-ahead dispatch.
\end{itemize}

\vspace{1ex}
\noindent\textbf{Project: Distributed Energy Management Frameworks} \hfill \textit{2024 -- 2025}
\begin{itemize}[leftmargin=0.25in, itemsep=0pt, topsep=0pt]
    \item Developed a 3-layer architecture linking \textbf{MILP} for discrete unit commitment, \textbf{ADMM} for decentralizing tie-lines without exposing private data, and \textbf{Nucleolus} for fair \textbf{Peer-to-Peer (P2P)} profit sharing.
    \item Applied \textbf{Demand Side Management (DSM)} to optimized local microgrid load scheduling for mitigating peak-load stress and reducing operational costs during islanded mode.
    \item Implemented an adaptive penalty tuning mechanism based on primal-dual residual balancing to resolve \textbf{ADMM} divergence caused by constrained P2P parameters, ensuring rapid convergence.
\end{itemize}

\vspace{1ex}
\noindent\textbf{Lab Contribution \& Mentorship} \hfill \textit{2024 -- Present}
\begin{itemize}[leftmargin=0.25in, itemsep=0pt, topsep=0pt]
    \item Architected custom Multi-Agent AI workflows to accelerate personal literature synthesis.
    \item Mentored junior researchers in academic reading methodologies and algorithmic modeling, driving lab research productivity.
\end{itemize}

\vspace{1ex}

% ACADEMIC PROJECTS
\section*{Academic Projects}

\noindent\textbf{Machine Learning for EV Load Forecasting} \hfill \textit{Mar. 2026 -- April 2026}
\begin{itemize}[leftmargin=0.25in, itemsep=0pt, topsep=0pt]
    \item Engineered a Proxy-Lag Cascade Ensemble architecture utilizing dual XGBoost models (t+24h and t+1) to eliminate phase lag and accurately forecast peak EV load demand.
\end{itemize}

\vspace{1ex}
\noindent\textbf{Industrial Power Supply \& Distribution Network Design} \hfill \textit{Oct. 2025 -- Dec. 2025}
\begin{itemize}[leftmargin=0.25in, itemsep=0pt, topsep=0pt]
    \item Designed a high/low-voltage distribution network for a large-scale manufacturing plant, optimizing the techno-economic balance between CAPEX and OPEX without compromising system stability.
    \item Sized transformers, optimized cable cross sections, and performed short circuit calculations to prevent thermal overloads and ensure reliability across all operating conditions.
\end{itemize}

\vspace{1ex}
\noindent\textbf{Electronic \& Hardware Prototyping} \hfill \textit{2023 -- 2024}
\begin{itemize}[leftmargin=0.25in, itemsep=0pt, topsep=0pt]
    \item Designed and physically prototyped an AC current measurement circuit and a DC-DC Buck converter to ensure stable voltage regulation and high-precision monitoring for hardware testing.
    \item Simulated layouts using LTspice, Proteus, and MATLAB/Simulink to validate circuit robustness prior to physical fabrication.
\end{itemize}

\vspace{1ex}

% PUBLICATIONS
\section*{Publications \& Presentations}

\begin{enumerate}[label={[\arabic*]}, leftmargin=0.25in, itemsep=0pt, topsep=0pt]
    \item \textbf{T. X. Thanh (1st Author)}, N. T. Anh, T. X. Hung, N. D. Tuyen, T. T. Son, G. Fujita, "A Decentralized MILP-ADMM Framework for P2P Energy Trading in Interconnected Microgrids," \textit{10th International Conference on Green Energy and Applications (ICGEA)}, Mar. 2026. \textbf{[Oral Presentation]}
    \item L. A. Quan, \textbf{T. X. Thanh (2nd Author)}, N. T. Anh, P. T. Vinh, T. X. Hung, N. D. Tuyen, T. T. Son, "Microgrid Optimization with MILP-based Demand Side Management," \textit{Journal of Science and Technology (JST)}, vol. 36, Jun. 2026. \textbf{[Oral Presentation at Vietnam Student Forum]}
\end{enumerate}

\vspace{1ex}

% SKILLS
\section*{Skills}

\begin{itemize}[leftmargin=0in, label={}, labelsep=0pt, itemsep=0pt, topsep=0pt]
    \item \textbf{Optimization:} MILP, SOCP, Convex Relaxation
    \item \textbf{Control \& Algorithms:} ADMM, ATC, MPC, Game Theory
    \item \textbf{Programming:} Python (Pyomo), C, C++
    \item \textbf{EE Simulation:} Simulink, LTspice, PVsyst, Proteus
    \item \textbf{Tools \& AI:} Multi-Agent AI, LaTeX, Git
    \item \textbf{Languages:} English (IELTS 6.0, Target 7.0)
\end{itemize}

\vspace{1ex}

% HONORS
\section*{Honors \& Awards}

\noindent 2026 \quad Best Presentation Award, 10th Int. Conf. on Green Energy and Applications (ICGEA)

\vspace{1ex}

% REFERENCES
\section*{References}

\noindent \textbf{Assoc. Prof. Nguyen Duc Tuyen} \\
\noindent Associate Professor and Head of PGRE Lab, Power Systems Department, HUST, Vietnam \\
\noindent Adjunct Assoc. Prof, Shibaura Institute of Technology, Japan \\
\noindent E-mail: \href{mailto:tuyen.nguyenduc@hust.edu.vn}{tuyen.nguyenduc@hust.edu.vn}

\end{document}
